\documentclass{article}
\usepackage[utf8]{inputenc}
\usepackage[T1]{fontenc}
\usepackage{amsmath,amssymb}
\usepackage{braket}

\begin{document}

\section*{Dos $n$‑qubits en estados base $\ket{x}, \ket{y}$}

Dos $n$‑qubits en estados base $\ket{x}, \ket{y}$ son ortogonales \emph{sii} $x \neq y$. Esto utilizando producto tensorial:
\[
\braket{x}{y}
= \braket{x_1}{y_1}\,\otimes\,\braket{x_2}{y_2}\,\otimes\cdots\otimes\braket{x_n}{y_n}.
\]
Tener en cuenta $\braket{x_k}{y_k} = \delta_{x_k y_k}$. Se tiene que los estados $\ket{x}$ y $\ket{y}$ son ortogonales si existe una posición $k$, en la cual los bits de $x$ y $y$ difieren.

\medskip
\textbf{Nota:}
\begin{gather*}
\ket{a}\otimes\ket{b} = \ket{ab},
\quad
\ket{c}\otimes\ket{d} = \ket{cd},\\
\braket{ab}{cd} = \braket{a}{c}\,\braket{b}{d}.
\end{gather*}

\section*{Ejemplo: Medición de un qubit}

Las mediciones sobre el primer qubit consistirán en realizar una medición proyectiva:
\[
P_j \;=\;\ket{j}\bra{j}\;\otimes\;I_{2^n}
\quad
\underbrace{\text{``Proyectores sobre el estado $j$''}}_{\text{}}.
\]
Para
\[
\ket{\psi}
= \sqrt{\tfrac{1}{3}}\,\ket{0}\ket{\phi}
\;+\;
\sqrt{\tfrac{2}{3}}\,\ket{1}\ket{\xi}.
\]
Si el resultado de la medición es $0$, el estado colapsa a $\ket{0}\ket{\phi}$ con una probabilidad de $\tfrac{1}{3}$. (Análogo si el resultado es $1$.)

\section*{Ejemplo: Ortogonalidad y registro cuántico}

\[
\braket{10}{11}
= \braket{1}{1}\,\braket{0}{1}
= 0.
\]

Un registro cuántico consistiría de un conjunto de $n$‑qubits, que, de estar en cualquier superposición,
\[
\alpha_0\ket{0} + \alpha_1\ket{1} + \dots + \alpha_N\ket{N},
\]
sujeto a la normalización
\[
\sum_{j=0}^{2^n-1} \lvert \alpha_j\rvert^2 = 1,
\qquad
N = 2^n - 1.
\]
Medir esto en la base computacional resultaría en un estado de $n$‑bits $\ket{j}$ con probabilidad $\lvert\alpha_j\rvert^2$.

\end{document}
